% ============================================================
% CAPÍTULO 1 - INTRODUÇÃO
% ============================================================

\chapter{INTRODUÇÃO}
\label{chap:introducao}

A avaliação da aprendizagem escolar constitui um dos temas mais complexos do campo educacional contemporâneo. Frequentemente reduzida a práticas burocráticas de classificação e seleção, a avaliação tem seu potencial formativo obscurecido por uma tradição histórica que a vinculou aos exames e à pedagogia do controle. Nesse contexto, a obra de Cipriano Carlos Luckesi, \textit{Avaliação da aprendizagem escolar: estudos e proposições}, destaca-se como um marco teórico fundamental para a compreensão crítica desse fenômeno e para a proposição de práticas avaliativas comprometidas com a democratização do ensino.

A presente leitura comentada analisa dois capítulos específicos: o Capítulo VII, \textit{Avaliação do aluno: a favor ou contra a democratização do ensino?}, e o Capítulo VIII, \textit{Planejamento e avaliação na escola: articulação e necessária determinação ideológica}. A escolha justifica-se pela relevância dos temas e pela articulação que estabelecem entre avaliação, democratização do ensino e planejamento pedagógico, elementos centrais para a compreensão da prática avaliativa em sua dimensão política e pedagógica. A abordagem de Luckesi caracteriza-se por uma dupla perspectiva: crítica, ao desvelar os determinantes histórico-sociais da prática avaliativa, e propositiva, ao anunciar novas possibilidades de ação \cite{luckesi-avaprend}.

A estrutura desta leitura comentada organiza-se em duas seções principais: a primeira analisa o Capítulo VII, examinando a definição de democratização do ensino, a crítica à prática atual de avaliação e a proposição da avaliação diagnóstica; a segunda analisa o Capítulo VIII, abordando a relação entre planejamento e avaliação e a dimensão ideológica do ato de planejar.

% ============================================================
% FIM DO ARQUIVO intro.tex
% ============================================================